% ===================================================================
% preamble.tex - Preambule du rapport "Retraite par les dividendes"
% Reprise de Business_Analysen/Neste/preamble.tex (tache 10, passage
% infrastructure): langue francaise, macros specifiques a Neste retirees
% (kurzGB*, tag*, tableaux et figure Neste), mecanique de citation
% generique conservee (\Q, \QL, \QR, \QI, \QW), macros de figure pgfplots
% et de citation par cle JSON (\QH, \QP) ajoutees pour ce projet.
% ===================================================================
\usepackage[T1]{fontenc}
\usepackage[utf8]{inputenc}
\usepackage{lmodern}
\usepackage[french]{babel}
% babel-french rend « : ; ! ? » actifs (espace fine automatique avant ces
% signes) - casse la syntaxe TikZ/pgfplots, qui utilise ";" comme
% terminateur de commande (\addplot ... ;) et ":" dans plusieurs cles.
% Desactive globalement (perte de l'espace fine automatique devant ces
% quatre signes en prose, negligeable ici). DOIT passer par
% \AtBeginDocument: un \shorthandoff{} place directement ici, avant
% \begin{document}, est silencieusement annule par le \selectlanguage{french}
% que babel declenche a \begin{document} (bascule de langue, pas au
% chargement du paquet) - piege confirme par un test de compilation reel de
% \linienfigur/\balkenfigur sur deux CSV (tache 10): un \shorthandoff{} nu
% ici laisse pdflatex echouer avec "Runaway argument? File ended while
% scanning use of \pgfplots@addplotimpl@table@fromfile" des le premier
% \addplot table{...} employe DANS une macro (\newcommand), alors que le
% meme code tape directement dans le corps du document compile; place apres
% \begin{document} (ou via \AtBeginDocument, meme effet, ordre des hooks),
% le \shorthandoff{} tient et les deux macros de figure compilent.
\AtBeginDocument{\shorthandoff{:;!?}}
\usepackage{csquotes}
% patch sans "footnote": microtype ne peut pas poser son patch de note de
% bas de page si footmisc (para) a deja remplace l'appareil de notes.
\usepackage[patch={item,toc,eqnum}]{microtype}
\usepackage{xcolor}
\usepackage{amsmath}
\usepackage{amssymb}
\usepackage{booktabs}
\usepackage{longtable}
\usepackage{array}
% \rowcolor[gray]{0.9} pour la ligne d'en-tete des tableaux (regle des flottants:
% premiere ligne grisee a 10 %). Charge apres array, dont il etend les colonnes.
\usepackage{colortbl}
\usepackage{tabularx}
\usepackage{siunitx}
\usepackage{graphicx}
\usepackage{tikz}
\usepackage{placeins}
\usepackage{needspace}
\usepackage{pgfplots}
\usepackage{pgfplotstable}
\usepackage[hidelinks]{hyperref}
\usepackage{url}
% xparse: argument optionnel des macros de figure (\linienfigur, \balkenfigur),
% correctif infrastructure tache 10 - voir plus bas.
\usepackage{xparse}

\areaset[current]{463.07pt}{636.6pt}

% Interdire veuves et orphelins (regle typographique heritee des rapports
% precedents, valable aussi en francais).
\clubpenalty=10000
\widowpenalty=10000
\displaywidowpenalty=10000

% Appareil de notes de bas de page a plat: toutes les notes d'une page dans
% un seul paragraphe continu. Le mode a deux colonnes (dblfnote) boucle avec
% pgfplots et hyperref dans la routine de sortie (heritee des rapports
% precedents).
\usepackage[para]{footmisc}

\usepackage[automark]{scrlayer-scrpage}
\pagestyle{scrheadings}
\renewcommand{\sectionmark}[1]{\markright{\thesection\quad #1}}
\clearpairofpagestyles
\setkomafont{pagehead}{\normalfont\small}
\ihead{\rightmark}
\ohead{Retraite par les dividendes}
\KOMAoptions{headsepline=true}
\cfoot*{\pagemark}

\setcounter{tocdepth}{2}

\pgfplotsset{compat=1.18}
% Pas de usepgfplotslibrary{dateplot}: tous les CSV de figures de ce projet
% portent une colonne "jahr" (annee entiere), jamais une date ISO; la
% bibliotheque dateplot n'a donc aucun usage ici (decision tache 10).
\pgfplotsset{/pgf/number format/use comma, /pgf/number format/1000 sep={\,}}
% Separateur de colonne par defaut pour toute lecture de table pgfplots
% (\addplot table[...] et \pgfplotstableread).
\pgfplotsset{table/col sep=comma}

% --- Format des nombres: virgule decimale francaise ----------------------
\sisetup{
  output-decimal-marker = {,},
  group-separator       = {\,},
  group-minimum-digits  = 4,
  % Sans group-digits=integer, siunitx groupe aussi les decimales:
  % 43.4679 deviendrait "43,467\,9" au lieu de "43,4679".
  group-digits          = integer,
  detect-weight         = true,
  detect-family         = true
}

% =====================================================================
% Figures pgfplots: style commun et deux macros generiques.
% pgfplots 1.18.1 est installee sur cette machine (verifie via
% kpsewhich/pgfplots.revision.tex). La cle "symbolic x coords from table"
% suggeree par le brief de la tache 10 N'EXISTE PAS dans cette version
% (absente de pgfplots.code.tex, verifie par grep; ni dans aucune version
% documentee de pgfplots - vraisemblablement une confusion avec
% "xticklabels from table", qui existe bel et bien). \balkenfigur utilise
% donc directement le REPLI que le brief prevoit pour ce cas: lire le CSV
% avec \pgfplotstableread, positionner les barres par index de ligne
% (x expr=\coordindex) et poser les etiquettes de graduation avec
% "xticklabels from table". Confirme par une compilation reelle sur
% data/fig_laender.csv (tache 10, voir aussi task-10-report.md).
% =====================================================================
\pgfplotsset{retraite/.style={
  width=\linewidth, height=6cm, grid=major, grid style={gray!20},
  tick label style={font=\footnotesize}, label style={font=\footnotesize},
  legend style={font=\footnotesize, draw=none, fill=none},
  /pgf/number format/use comma, /pgf/number format/1000 sep={\,}}}

% \linienfigur{fichier}[cles pgfplots]{colonne x}{liste colonnes y}{liste legende}{caption+label}
% fichier: nom du CSV dans data/ (sans extension). colonne x: nom de colonne.
% liste colonnes y et liste legende: listes separees par des virgules, dans
% le MEME ORDRE (ex. {brutto,netto}{Brut,Net}). caption+label: texte de
% \caption, \label{fig:...} a la fin (convention code-style.md), ex.
% "Salaire brut et net simule au fil des annees.\label{fig:salaire}".
%
% [cles pgfplots] (correctif infrastructure, tache 10, ARGUMENT OPTIONNEL,
% retro-compatible: un appel a 5 arguments obligatoires comme avant, sans
% [...], continue de compiler a l'identique): liste de cles pgfplots
% standard, ex. [xlabel={Ann\'ee}, ylabel={Cible mensuelle (euros nominaux)}].
% Appliquees APRES le style "retraite" et apres le xlabel par defaut
% (#2, nom brut de colonne), donc une cle repetee ici l'emporte - remplace
% le contournement "{\pgfplotsset{every axis post/.append style={...}}}
% \linienfigur{...}}" utilise par les chapitres 2-3 avant ce correctif.
% Place APRES {fichier} (pas avant, malgre la convention LaTeX habituelle du
% "\sqrt[n]{x}"): {fichier} reste ainsi colle au nom de macro sur la MEME
% ligne source meme quand le bloc d'options s'etend sur plusieurs lignes,
% ce dont depend scripts/check_literals.py (l'exemption du nom de CSV, qui
% contient souvent un nombre, ex. fig_x_250, marche ligne par ligne).
\NewDocumentCommand{\linienfigur}{m O{} m m m m}{%
\begin{figure}[htbp]\centering\begin{tikzpicture}\begin{axis}[retraite, xlabel={#3}, #2]
\foreach \y in {#4}{\addplot+[thick, mark=none] table[x=#3, y=\y, col sep=comma]{data/#1.csv};}
\legend{#5}\end{axis}\end{tikzpicture}\caption{#6}\end{figure}}

% \balkenfigur{fichier}[cles pgfplots]{colonne x (etiquettes)}{colonne y}{caption+label}
% fichier/colonne x/colonne y/caption+label: memes conventions que
% \linienfigur. Barres positionnees par index de ligne (x expr=\coordindex,
% axe numerique). [cles pgfplots]: meme argument optionnel que \linienfigur
% (ex. [ylabel={Montant annuel (euros de 2026)}, bar width=22pt]), place
% APRES {fichier} pour la meme raison (voir commentaire de \linienfigur).
%
% Etiquettes de graduation (correctif infrastructure, tache 10): posees par
% "xticklabels from table" (repli du brief, la cle "symbolic x coords from
% table" n'existant pas dans pgfplots 1.18.1, voir commentaire au-dessus) sur
% la colonne "#3_disp", PAS #3 directement. #3_disp est ecrite par
% scripts/rechnung_retraite.py (schreiben()) pour toute colonne texte: son
% contenu est echappe pour LaTeX (escapieren()), de sorte qu'une valeur de
% colonne contenant un underscore (ex. "brut_necessaire", "capital_10000")
% compile toujours SANS INTERVENTION DE LA REDACTION (l'underscore, catcode 8
% hors mode mathematique, ferait planter LaTeX si on le lisait tel quel).
% Pour des etiquettes plus lisibles que la colonne brute echappee (ex.
% "Dividende brut" plutot que "brut\_necessaire"), passer une cle
% "xticklabels={...}" explicite dans [cles pgfplots]: appliquee APRES la
% valeur par defaut ci-dessous, elle la remplace (meme mecanique que pour
% \linienfigur) - c'est la voie normale pour les etiquettes en prose.
\NewDocumentCommand{\balkenfigur}{m O{} m m m}{%
\begin{figure}[htbp]\centering\begin{tikzpicture}
\pgfplotstableread{data/#1.csv}\balkendaten
\begin{axis}[retraite, ybar, xtick=data,
  xticklabels from table={\balkendaten}{#3_disp},
  x tick label style={rotate=35, anchor=east}, #2]
\addplot table[x expr=\coordindex, y=#4]{\balkendaten};\end{axis}\end{tikzpicture}
\caption{#5}\end{figure}}

% --- Citation de documents (mecanique generique, conservee de Neste) -----
% \Q{Dokument}{Seite}: note avec titre court (saut vers le repertoire des
% sources) et page IMPRIMEE. Le titre court \kurz<Dokument> doit etre
% defini par qui appelle \Q (aucun \kurz* predefini ici: ce projet source
% l'essentiel de ses valeurs via data/quellen.json et \QH, voir plus bas;
% \Q reste disponible pour une citation ponctuelle d'un document de refs/
% sans entree quellen.json, par exemple le classeur Shiller).
% \QL{Schluessel}{Dokument}{Seite}: note avec marque; \QR{Schluessel}: la
% MEME note re-referencee (seulement si meme page imprimee, verifie par
% scripts/check_footnote_groups.py). \QI: reference dans le texte, pour une
% cellule de tableau ou une note de bas de page se perdrait dans un flottant.
\definecolor{quellenlink}{RGB}{20,60,120}
\newcommand{\quellensprung}[2]{\hyperref[#1]{\textcolor{quellenlink}{#2}}}
\newcommand{\Q}[2]{\footnote{\quellensprung{quelle:#1}{\csname kurz#1\endcsname}, p.~#2.}}
\newcommand{\QL}[3]{\footnote{\label{fn:#1}\quellensprung{quelle:#2}{\csname kurz#2\endcsname}, p.~#3.}}
\newcommand{\QR}[1]{\textsuperscript{\ref{fn:#1}}}
\newcommand{\QI}[2]{{\scriptsize[\csname tag#1\endcsname\ p.\,#2]}}
\newcommand{\QW}[3]{\footnote{\quellensprung{#2}{#1} (consult\'e le #3).}}

% --- Citation par cle JSON (data/quellen.json, data/presse.json) ---------
% \QH{cle}: note de bas de page complete (titre/organisme, reference ou
% page si presente, URL, date de consultation) pour une entree de
% data/quellen.json. Generee par scripts/rechnung_retraite.py
% (quellen_tex()) dans data/quellen.tex (macro \QHfn@<cle>, \input dans
% analyse.tex) et data/quellen_annexe.tex (liste labellisee
% \label{src:<cle>}, destinee a l'annexe "Sources"). Usage:
%   Le taux reduit est de \KvSatzReduit~\%\QH{kv_satz_ermaessigt}.
% \QHtexte{cle} donne seulement le texte de la note (sans \footnote), pour
% ancrer soi-meme une citation qui revient sur la meme page imprimee:
%   \footnote{\label{fn:monancre}\QHtexte{cle}} ... plus loin ... \QR{monancre}
% (verifie alors par scripts/check_footnote_groups.py comme un \QL/\QR).
\newcommand{\QHtexte}[1]{\csname QHfn@#1\endcsname}
\newcommand{\QH}[1]{\footnote{\QHtexte{#1}}}
% \QP{id}: meme mecanique pour une source de presse (data/presse.json),
% dedupliquee par id de source (plusieurs citations du tableau
% data/tab_presse.tex, chapitre 9, peuvent partager le meme id). Le corps
% de note reprend quelle/titel/url/abgerufen, jamais la citation elle-meme
% (celle-ci reste dans le tableau, pas dans la note).
\newcommand{\QPtexte}[1]{\csname QPfn@#1\endcsname}
\newcommand{\QP}[1]{\footnote{\QPtexte{#1}}}

% --- Unites: les macros de data/kennzahlen.tex sont deja mises en forme --
\newcommand{\Mio}[2]{#1\,Mio.\,#2}
\newcommand{\je}[2]{#1\,#2}
\newcommand{\Pz}[1]{#1\,\%}

% --- Trois couleurs, trois responsabilites --------------------------------
% Noir: faits et calculs sources. Bleu: appreciation du redacteur.
% Rouge: verdict/recommandation.
\definecolor{vdrot}{RGB}{155,25,30}
\newcommand{\vd}[1]{{\color{vdrot}#1}}
\definecolor{bkblau}{RGB}{20,60,150}
\newcommand{\bk}[1]{{\color{bkblau}#1}}

% --- Corps de tableau genere, sans \relax accroche ------------------------
\makeatletter
\newcommand{\tabellenkoerper}[1]{\@@input #1.tex }
\makeatother

\newcommand{\annahme}[1]{%
  \par\medskip\noindent\fbox{\parbox{0.97\linewidth}{\small\textbf{Hypoth\`ese~:} #1}}\par\medskip}

% --- Encadre "Yann en <annee>": un exemple chiffre par chapitre ----------
% \Yannbox{annee}{texte}: annee est TOUJOURS une macro (ex. \AnneeYannBasis),
% jamais un chiffre tape a la main (regle du depot, verifiee par
% scripts/check_literals.py, les annees restant autorisees telles quelles);
% texte ne doit lui aussi contenir que des macros de data/kennzahlen.tex et
% des annees.
\definecolor{yanncadre}{RGB}{20,110,80}
\newcommand{\Yannbox}[2]{%
  \par\medskip\noindent
  {\color{yanncadre}\rule{0.97\linewidth}{0.6pt}}\par\nobreak
  \noindent\parbox{0.97\linewidth}{\small\textbf{\color{yanncadre}Yann en #1.} #2}%
  \par\nobreak\noindent{\color{yanncadre}\rule{0.97\linewidth}{0.6pt}}\par\medskip}
